\chapter{Metodología}\label{ch:metodologia}

% ============================================================
%  VISIÓN GENERAL DEL CAPÍTULO
% ============================================================

\section{Diseño general}\label{sec:diseno_general}

El objetivo metodológico de esta tesis consiste en caracterizar la distribución espacial y la magnitud de los máximos de precipitación y viento asociados a los ciclones extratropicales del Atlántico Sur, adoptando un marco de referencia lagrangiano centrado en el sistema. El análisis se condiciona sistemáticamente por dos factores: la fase del ciclo de vida y la región de ciclogénesis. La estrategia experimental se estructura en una secuencia lógica de cinco etapas: (1) definición de las fuentes de datos; (2) construcción y estratificación de las poblaciones de ciclones; (3) procesamiento lagrangiano de los campos de superficie; (4) caracterización estadística univariada mediante estimación de densidad y análisis EOF; y (5) análisis multivariado (MEOF) y construcción de prototipos estructurales. Cada etapa produce el insumo de la siguiente, configurando un pipeline reproducible y auditable.

% ============================================================
%  SECCIÓN M1 — DATOS
% ============================================================

\section{Datos}\label{sec:datos}

\subsection{Reanálisis ERA5}\label{sec:datos_era5}

Para la caracterización de los campos atmosféricos se utilizó el reanálisis de quinta generación del ECMWF, ERA5 \citep{hersbach2020era5}, accediendo específicamente al conjunto de datos \textit{``ERA5 hourly data on single levels from 1940 to present''} \citep{hersbach_era5_2023b} para el período 1979--2020. La justificación de la elección de ERA5 ---incluyendo su resolución espacial ($\sim$31 km) y temporal (horaria) y su representación superior de la intensidad ciclónica en el Atlántico Sur frente a generaciones previas de reanálisis \citep{gramcianinov2020analysis}--- fue presentada en la Introducción (Sección~\ref{sec:definicion_regiones}) y en los Fundamentos (Sección~\ref{sec:tb_fundamentos}). No obstante, debe reconocerse una limitación documentada en la representación de extremos: \citet{chen2024evaluation} demuestran que ERA5 subestima sistemáticamente los extremos del 5\% de ciclones extratropicales más intensos, con sesgos del $-10.2\%$ en viento y $-22.6\%$ en precipitación respecto a observaciones en Estados Unidos de America. Esta tendencia implica que los resultados obtenidos para el subconjunto $p90$ deben interpretarse como estimaciones conservadoras de la severidad real de los eventos. Esta sección se limita a detallar las variables de superficie seleccionadas y los parámetros de extracción.

% Para la caracterización de los campos atmosféricos se utilizó el reanálisis de quinta generación del ECMWF, ERA5 \citep{hersbach2020era5}, accediendo específicamente al conjunto de datos \textit{``ERA5 hourly data on single levels from 1940 to present''} \citep{hersbach_era5_2023b} para el período 1979--2020. La justificación de la elección de ERA5 —incluyendo su resolución espacial ($\sim$31 km) y temporal (horaria), su representación en el Atlántico Sur frente a generaciones previas de reanálisis, y sus limitaciones en la cuantificación— fue presentada en la Introducción y en los Fundamentos (Secciones~\ref{sec:tb_fundamentos} y~\ref{subsec:tb_compound_def}). Esta sección se limita a detallar las variables de superficie seleccionadas y los parámetros de extracción.

\subsection{Variables meteorológicas}\label{subsec:variables}

La caracterización de los patrones espaciales se realizó a partir de campos horarios extraídos de ERA5, utilizando las siguientes variables de superficie:

\begin{itemize}
    \item Precipitación total acumulada (\texttt{tp}, en mm\,h$^{-1}$).
    \item Magnitud del viento en superficie a 10 m (\texttt{wind10}, en m\,s$^{-1}$).
\end{itemize}

La caracterización conjunta de ambas variables responde a la naturaleza morfológica de los ciclones extratropicales y su respectiva estructura de los campos de precipitación y viento, fundamentado en la Sección~\ref{subsec:tb_compound_def} de los Fundamentos . El análisis simultáneo permite describir la configuración espacial típica de los campos de superficie en cada fase evolutiva, sirviendo además como referencia para evaluar, la eventual co-ocurrencia de los máximos de ambas variables. Esta estrategia informa la posterior aplicación del análisis multivariado (MEOF) detallado en la Sección~\ref{subsec:eof_multivariado}.

\subsection{Base de trayectorias y fases}\label{sec:bd_ciclones}

Para la identificación y caracterización de los sistemas, esta investigación utiliza la base de datos procesada por \citet{coutodesouza2024new}, quien tomó el catálogo original de trayectorias desarrollado por \citet{gramcianinov2019properties} y le incorporó una clasificación objetiva de las fases evolutivas mediante el algoritmo \textit{CycloPhaser} \citep{deSouza2025CycloPhaser}. Los fundamentos físicos que justifican el uso de la vorticidad relativa ($\zeta_{850}$) como variable de tracking y la segmentación objetiva del ciclo de vida en cuatro fases fueron expuestos en las Secciones~\ref{subsec:tb_zeta} y~\ref{subsec:tb_fases} de los Fundamentos. Esta sección se limita a describir las especificaciones técnicas y los criterios de filtrado del conjunto de datos.

En su construcción original \citep{gramcianinov2019properties}, el catálogo se deriva de los campos del reanálisis ERA5, aprovechando su alta resolución para capturar sistemas de mesoescala y ciclogénesis orográficas comunes en la región de estudio \citep{gramcianinov2020analysis}. Los parámetros de configuración del algoritmo y los filtros aplicados se detallan en la Tabla~\ref{tab:tracking_params}.

\begin{table}[ht]
\centering
\caption[Parámetros y filtros de la base de trayectorias]{Parámetros técnicos y filtros aplicados en la base de datos de trayectorias \citep{gramcianinov2020analysis}.}
\label{tab:tracking_params}
\begin{tabular*}{\textwidth}{@{\extracolsep{\fill}}ll}
\hline
Parámetro & Especificación \\ 
\hline
Fuente de datos & ERA5 (ECMWF) \\
Resolución espacial & $\sim$31 km ($0.25^\circ \times 0.25^\circ$) \\
Variable de seguimiento & Vorticidad relativa en 850 hPa ($\zeta_{850}$) \\
% % Nivel de filtrado espectral & T42 (para aislar la escala sinóptica) \\
% % Umbral de intensidad & $\zeta_{850} < -1.0 \times 10^{-5}$ s$^{-1}$ \\
Duración mínima & 24 horas \\
Desplazamiento mínimo & 1000 km \\ 
\hline
\end{tabular*}
\end{table}

% \begin{table}[ht]
% \centering
% \caption{Parámetros técnicos y filtros aplicados en la base de datos de trayectorias \citep{gramcianinov2019properties}.}
% \label{tab:tracking_params}
% \begin{tabular*}{\textwidth}{@{\extracolsep{\fill}}ll}
% \hline
% Parámetro & Especificación \\ 
% \hline
% Fuente de datos & ERA5 (ECMWF) \\
% Resolución espacial & $\sim$31 km ($0.25^\circ \times 0.25^\circ$) \\
% Variable de seguimiento & Vorticidad relativa en 850 hPa ($\zeta_{850}$) \\
% % % Nivel de filtrado espectral & T42 (para aislar la escala sinóptica) \\
% % % Umbral de intensidad & $\zeta_{850} < -1.0 \times 10^{-5}$ s$^{-1}$ \\
% Duración mínima & 24 horas \\
% Desplazamiento mínimo & 1000 km \\ 
% \hline
% \end{tabular*}
% \end{table}


Sobre este seguimiento espacial continuo, \citet{coutodesouza2024new} asignan una etiqueta de fase a cada paso de tiempo de la trayectoria mediante \textit{CycloPhaser}, evaluando la tasa de cambio de la vorticidad central para aislar con precisión matemática los momentos de desarrollo y decaimiento del sistema. La Tabla~\ref{tab:fases_def} resume las cuatro fases evolutivas consideradas en este estudio.

\input{tabelas_es/tab_fases_def.tex}

% ============================================================
%  SECCIÓN M2 — POBLACIONES DE ESTUDIO
% ============================================================

\section{Construcción de las poblaciones de estudio}\label{sec:poblaciones_estudio}

La segmentación de la muestra sigue los criterios conceptuales presentados en la Sección~\ref{subsec:tb_poblaciones} de los Fundamentos. La base de datos de trayectorias fue dividida en dos niveles jerárquicos: una población climatológica de referencia (Población Global) y, a partir de ella, un subconjunto de alta intensidad (p90). Ambas poblaciones se construyeron de forma independiente para cada una de las tres regiones de ciclogénesis (SBR, LPB y ARG), cuya delimitación espacial sigue la taxonomía propuesta por \citet{gramcianinov2019properties} y se detalla en la Tabla~\ref{tab:regiones_coords}. La estratificación geográfica responde a que estas regiones presentan regímenes dinámicos diferenciados —subtropical/híbrido para SBR, transicional para LPB y baroclínico para ARG—, cuyas características físicas fueron discutidas en la Sección~\ref{sec:tb_regiones} de los Fundamentos.

\input{tabelas_es/tab_regiones_coords.tex}

\subsection{Población Global: climatología de referencia}\label{subsec:grupo_global}

El objetivo de la Población Global es establecer el comportamiento morfológico y estadístico medio de los ciclones con un desarrollo típico y completo en la región. Para conformarla, se aplicó un filtro de coherencia evolutiva a la base de ciclones, seleccionando únicamente los sistemas que cumplieron con dos criterios simultáneos:

\begin{enumerate}
    \item Ciclo de vida completo: Presencia confirmada y secuencial de las cuatro fases de desarrollo —incipiente (\textit{incipient}), intensificación (\textit{intensification}), madurez (\textit{mature}) y decaimiento (\textit{decay})—.
    \item Ausencia de fases complejas: Exclusión categórica de sistemas que desarrollaron fases secundarias, de re-intensificación (e.g., \textit{intensification~2}, \textit{mature~2}) o fases residuales (\textit{residual}).
\end{enumerate}

Este doble filtrado elimina el ruido introducido por vórtices efímeros o de evolución errática, garantizando una muestra climatológica homogénea. Para cada ciclón de la Población Global se identificó el registro temporal correspondiente a su máxima intensidad —el extremo absoluto de vorticidad relativa a 850 hPa, $\zeta_{850}^{\min}$— como métrica de intensidad para la segmentación posterior.

\subsection{Subconjunto de alta intensidad (p90)}\label{subsec:subconjuntos_intensidad}

Para evaluar la sensibilidad de los patrones espaciales frente a la severidad dinámica del sistema, se extrajo de la Población Global el 10\% de ciclones con los valores más negativos de $\zeta_{850}^{\min}$, correspondientes al primer decil de la distribución climatológica. Este estrato se denominará estrictamente subconjunto p90, siguiendo la definición conceptual y el respaldo metodológico expuestos en la Sección~\ref{subsec:tb_poblaciones} de los Fundamentos.

La Tabla~\ref{tab:conteo_ciclones} documenta la reducción del tamaño muestral desde la base cruda hasta la conformación de la Población Global y el subconjunto p90.

\input{tabelas_es/tab_conteo_ciclones.tex}

Un aspecto necesario de la estratificación es verificar que los sistemas del subconjunto p90 mantengan una evolución temporal coherente con el paradigma dinámico esperado. El análisis de la fase en que se registra el pico de vorticidad (Tabla~\ref{tab:fase_max_intensidad}) confirma que el subconjunto p90 respeta la estructura dinámica dominante: aproximadamente el 83\% de los casos alcanza su máxima intensidad durante la fase de madurez, frente a menos del 5\% durante la fase de decaimiento. Este comportamiento es coherente con el observado en la Población Global y respalda el criterio estadístico del percentil 90 para capturar sistemas maduros y organizados, excluyendo vórtices transitorios o de desarrollo difuso.

\input{tabelas_es/tab_fase_max_intensidad.tex}

\subsection{Análisis comparativo}\label{subsec:estrategia_analise}

Sobre el subconjunto p90 —y comparándolo sistemáticamente con la Población Global— se aplica íntegramente el flujo metodológico descrito en las secciones siguientes. Esta estrategia comparativa permitirá evaluar si una mayor intensidad dinámica del vórtice se asocia con cambios significativos en: (i) la magnitud absoluta de los máximos de viento y precipitación; (ii) la distribución espacial relativa de dichos máximos; y (iii) los patrones estructurales dominantes identificados por el MEOF.

% ============================================================
%  SECCIÓN M3 — PROCESAMIENTO LAGRANGIANO
% ============================================================

\section{Procesamiento lagrangiano}\label{sec:procesamiento_lagrangiano}

Para caracterizar la estructura de los sistemas independientemente de su ubicación geográfica, se adopta un marco de referencia lagrangiano expuesto en la Sección~\ref{subsec:tb_lagrangiano} de los Fundamentos. Se define un dominio móvil de $20^{\circ} \times 20^{\circ}$ centrado en la posición del máximo de vorticidad para cada instante $t$ de la trayectoria, en coordenadas relativas $(\Delta x, \Delta y)$ respecto al núcleo ciclónico. Este dominio permite aislar la señal de mesoescala y representar físicamente la extensión espacial de los brazos frontales del sistema.

El procesamiento se realiza de forma separada para cada combinación de región de ciclogénesis y fase del ciclo de vida, permitiendo estratificar los resultados simultáneamente por ambos factores.

Siguiendo la técnica de compositing centrado en el sistema descrita en la Sección~\ref{subsubsec:tb_compositing}, y con el objetivo de reducir la variabilidad de alta frecuencia inherente a los datos horarios e aislar la condición estructural más representativa de cada fase, se aplica un promediado centrado. Para cada ciclón y fase, se seleccionan los registros horarios correspondientes al núcleo temporal de dicha fase —los dos o tres pasos centrales según la duración total de la fase— y se promedian sus campos espaciales en coordenadas lagrangianas.

Este criterio garantiza que el campo resultante refleje la condición arquetípica del régimen (incipiente, intensificación, madurez o decaimiento) y no una mezcla de dos estados sucesivos, generando un único campo promedio por ciclón--fase que constituye la unidad de análisis de todas las etapas estadísticas subsiguientes.

% ============================================================
%  SECCIÓN M4 — CARACTERIZACIÓN ESTADÍSTICA UNIVARIADA
% ============================================================

\section{Caracterización estadística univariada}\label{sec:tecnicas_estadisticas}

\subsection{Extracción de máximos y distribución de probabilidad (PDF)}\label{subsec:extraccion_pdf}

A partir del campo promedio por ciclón--fase, se delimita el área de análisis aplicando una máscara circular de $9.5^{\circ}$ de radio centrada en el núcleo del sistema. Previo a la extracción de los máximos, se aplica un suavizado espacial bidimensional con un filtro gaussiano ($\sigma = 0.25$) mediante la biblioteca \texttt{scipy}, con el objetivo de mitigar el impacto de anomalías aisladas a escala de punto de grilla. Dentro de este radio de influencia, se identifican: (i) la magnitud del valor máximo absoluto de la variable de interés, $x_i$, y (ii) su ubicación en coordenadas relativas al centro del ciclón, $\mathbf{s}_i = (\Delta x_i, \Delta y_i)$. Este procedimiento genera, para cada combinación de región, fase evolutiva y variable, un conjunto de $m$ pares magnitud--ubicación (uno por evento ciclónico).

Para caracterizar la distribución estadística de las magnitudes máximas extraídas $\{x_i\}_{i=1}^{m}$, se estima la función de densidad de probabilidad (PDF) mediante KDE. La justificación de esta técnica frente a histogramas clásicos, así como su formulación matemática, se presentan en la Sección~\ref{subsubsec:tb_kde} de los Fundamentos. La implementación computacional utiliza la clase \texttt{gaussian\_kde} de la biblioteca \texttt{scipy.stats}, con el ancho de banda $h$ determinado automáticamente mediante la regla de Scott \ref{eq:tb_kde_pdf} y el núcleo gaussiano definido en la Ecuación~\ref{eq:tb_gaussian_kernel} de los Fundamentos. Las curvas PDF resultantes permiten contrastar cuantitativamente cómo evolucionan las magnitudes de las variables a lo largo del ciclo de vida y cómo difieren entre la Población Global y el subconjunto p90.

\subsection{Distribución espacial de los máximos}\label{subsec:kde_localizacion}

Para caracterizar la ubicación preferencial de los máximos respecto al centro del ciclón, se trata la posición relativa de cada máximo como una variable aleatoria bidimensional $\mathbf{S} = (\Delta x, \Delta y)$. A partir del conjunto de $m$ ubicaciones relativas $\{\mathbf{s}_i\}_{i=1}^{m}$ extraídas en la etapa anterior, se estima la función de densidad de probabilidad de la posición (PDFe), $\hat{f}(\mathbf{s})$, mediante el estimador kernel bidimensional presentado en la Sección~\ref{subsubsec:tb_kde} (Ecuaciones~\ref{eq:tb_kde_simple}--\ref{eq:tb_gauss_kernel_2d_simple}).

A diferencia de la PDF de magnitudes de la Sección~\ref{subsec:extraccion_pdf}, que describe la distribución de probabilidad de las \emph{magnitudes} de la variable, la PDFe cuantifica la densidad de probabilidad de que el máximo ocurra en una posición relativa dada $\mathbf{s} = (\Delta x, \Delta y)$, independientemente del valor numérico que tome la precipitación o el viento en ese punto. El resultado es un campo continuo de densidad de probabilidad sobre el dominio centrado.

Para su visualización se extraen dos niveles de contorno correspondientes a los percentiles 75 y 90 de la distribución acumulada de la densidad, calculados según:

\begin{equation}
    \hat{f}^*_q = \hat{f}_{(k)}, \quad \text{donde } k = \min\!\left\{j : \frac{\sum_{i=1}^{j}\hat{f}_{(i)}}{\sum_{i=1}^{m}\hat{f}_{(i)}} \geq q \right\},
\label{eq:nivel_contorno_kde}
\end{equation}

siendo $\hat{f}_{(i)}$ los valores de densidad ordenados ascendentemente y $q$ la fracción acumulada deseada ($q = 0.75$ ó $0.90$). La isolínea del percentil 75 delimita la región que concentra el 75\% de la masa de probabilidad espacial; la del percentil 90 encierra el núcleo más denso de ocurrencia de máximos. Adicionalmente, se identifica la coordenada modal:
\begin{equation}
    (\Delta x_{\text{moda}}, \Delta y_{\text{moda}}) = \underset{(\Delta x, \Delta y)}{\arg\max}\; \hat{f}(\Delta x, \Delta y),
\label{eq:moda_kde}
\end{equation}

que representa la posición relativa al centro del ciclón donde estadísticamente es más frecuente encontrar el máximo de la variable.

\subsection{Análisis de Funciones Ortogonales Empíricas (EOF) univariado}\label{subsec:eof_univariado}

Para identificar los patrones espaciales dominantes de variabilidad estructural, se aplica un análisis EOF sobre el conjunto de $m$ campos promedio por ciclón--fase, de forma separada para cada variable (precipitación, viento), región de ciclogénesis y fase evolutiva. Las referencias de investigaciones previas con este enfoque de descomposición modal fueron presentados en la Sección~\ref{subsubsec:tb_eof}. Previo al análisis, a cada campo se le resta la media del ensamble en cada punto de grilla para obtener las anomalías estructurales:

\begin{equation}
    X'(\tau, \mathbf{s}) = X(\tau, \mathbf{s}) - \frac{1}{m}\sum_{\tau=1}^{m}X(\tau, \mathbf{s}),
\end{equation}

donde $X(\tau, \mathbf{s})$ es el campo de la variable para el ciclón $\tau$ en la posición relativa $\mathbf{s} = (x,y)$. A partir de la matriz de anomalías se calcula la matriz de covarianza espacial, cuya diagonalización produce los modos ortogonales ($\text{EOF}_k$) y sus amplitudes ($\text{PC}_k$):

\begin{equation}
    X'(\tau, \mathbf{s}) = \sum_{k=1}^{K} \text{PC}_k(\tau) \cdot \text{EOF}_k(\mathbf{s}).
\end{equation}

No se aplica remoción de tendencia temporal (\textit{detrend}): el eje muestral está constituido por un índice de eventos discretos con espaciado temporal irregular a lo largo de 1979--2020, por lo que un ajuste lineal cronológico estándar asumiría erróneamente un espaciado constante. Desde el punto de vista físico, el objetivo es aislar las características morfológicas de la variabilidad inter-sistémica dentro de cada fase, no las tendencias climáticas de baja frecuencia. La implementación computacional utiliza la biblioteca \texttt{xeofs} \citep{rieger2024xeofs} sobre estructuras \texttt{xarray} \citep{hoyer2017xarray}.

\subsection{Significancia espacial: Bootstrap-t}\label{subsec:bootstrap_significancia}

Para evaluar la robustez estadística de los patrones espaciales (autovectores) derivados del EOF y aislar las señales físicas del ruido estocástico, se implementa el método de remuestreo no paramétrico Bootstrap-t. La elección de esta variante, así como sus ventajas frente al Bootstrap percentílico estándar para distribuciones asimétricas y de cola pesada, fueron fundamentadas en la Sección~\ref{subsubsec:tb_bootstrap} de los Fundamentos. Se generan $R = 15{,}000$ submuestras con reemplazo, garantizando que el error de Monte Carlo sea insignificante en la estimación de las colas de la distribución.

Para cada iteración $b$, se recalcula el análisis EOF sobre la submuestra y se obtiene el estadístico cuasi-pivotal:

\begin{equation}
    t^*_b = \frac{\hat{\theta}^*_b - \hat{\theta}}{SE^*_b},
\end{equation}

donde $\hat{\theta}$ es la carga espacial original, $\hat{\theta}^*_b$ su estimación en la submuestra y $SE^*_b$ su error estándar estimado. Los intervalos de confianza al 95\% se obtienen invirtiendo los cuantiles empíricos $q_{\alpha/2}$ y $q_{1-\alpha/2}$ de la distribución de $t^*$ (con $\alpha = 0.05$):

\begin{equation}
    \left(\hat{\theta} - q_{1-\alpha/2}\,SE,\; \hat{\theta} - q_{\alpha/2}\,SE\right).
\end{equation}

Los puntos de grilla se consideran estadísticamente significativos si su intervalo de confianza excluye el cero.

\subsection{Análisis de las Componentes Principales (PC1)}\label{subsec:analisis_pc}

Además de la descomposición espacial, se analizan las propiedades estadísticas de la Componente Principal del primer modo (PC1), que representa la amplitud de cada ciclón sobre el patrón estructural dominante. Este análisis permite inferir propiedades de la población que no son evidentes desde el patrón espacial. Se calculan para la serie de PC1 de cada variable, región y fase: la asimetría ($\gamma_1$):
\begin{equation}
    \gamma_1 = \frac{1}{m} \sum_{\tau=1}^{m} \left(\frac{PC_{1,\tau} - \overline{PC_1}}{\sigma_{PC_1}}\right)^3,
\end{equation}
y la curtosis ($\gamma_2$):
\begin{equation}
    \gamma_2 = \frac{1}{m} \sum_{\tau=1}^{m} \left(\frac{PC_{1,\tau} - \overline{PC_1}}{\sigma_{PC_1}}\right)^4 - 3.
\end{equation}
La asimetría positiva indica una subpoblación de ciclones con amplitud anormalmente alta en el modo dominante; la leptocurtosis ($\gamma_2 > 0$) sugiere colas pesadas y concentración central pronunciada, indicando que una subpoblación de sistemas se ajusta estrechamente al patrón arquetípico con estructuras intensas o atípicas.

Para establecer la conexión física entre la estructura de los campos de superficie y la intensidad del sistema, se evalúa la correlación de Pearson entre la PC1 de cada variable y el mínimo de vorticidad relativa alcanzado por cada ciclón, $\zeta_{850}^{\min}$:
\begin{equation}
    \rho\!\left(PC_1, \zeta_{850}^{\min}\right) = \frac{\text{cov}\!\left(PC_1, \zeta_{850}^{\min}\right)}{\sigma_{PC_1} \cdot \sigma_{\zeta}}.
\end{equation}
Una correlación significativa y negativa (dado que $\zeta_{850}^{\min}$ es negativa en el Hemisferio Sur) indicaría que los ciclones más intensos presentan mayor proyección sobre el patrón estructural dominante; una correlación débil sugeriría que la organización espacial de los impactos está modulada por forzantes locales independientes de la vorticidad del núcleo.

Para evaluar si los patrones dominantes de precipitación y viento varían conjuntamente, se calcula la correlación entre las PC1 de precipitación y viento dentro de cada región y fase:
\begin{equation}
    \rho = \text{corr}\!\left(PC_1^{\text{precip}},\; PC_1^{\text{viento}}\right).
\end{equation}
Valores cercanos a $+1$ indican que la precipitación y el viento se organizan espacialmente de manera coherente; valores cercanos a $0$ sugieren que la variabilidad espacial de la precipitación está controlada por procesos locales independientes de la configuración del campo de viento de gran escala.

% ============================================================
%  SECCIÓN M5 — ANÁLISIS MULTIVARIADO Y PROTOTIPOS
% ============================================================

\section{Análisis multivariado y prototipos estructurales}\label{sec:meof_prototipos}

\subsection{EOF combinado (MEOF) precipitación--viento}\label{subsec:eof_multivariado}

Para identificar los patrones espaciales dominantes de la estructura combinada de precipitación y viento, se realiza un análisis EOF combinado (MEOF). A diferencia del análisis univariado, el MEOF captura modos de variabilidad conjunta que reflejan la organización morfológica típica de los campos de superficie asociados a la dinámica ciclónica. Los fundamentos de esta técnica, incluyendo la necesidad de normalizar variables con unidades físicas dispares y la lógica de la concatenación espacial del vector de estado, fueron expuestos en la Sección~\ref{subsubsec:tb_eof} de los Fundamentos. Aunque el MEOF constituye una técnica habitual en estudios de eventos compuestos, en este trabajo se emplea exclusivamente como herramienta de síntesis estructural: su objetivo es aislar las configuraciones espaciales más recurrentes de los campos de superficie, no cuantificar la probabilidad de impactos múltiples. El procedimiento computacional consta de tres pasos:

\paragraph{Paso 1 — Anomalías.} Se calcula el campo de anomalías de cada variable respecto a la media de la Población Global en cada punto de grilla.

\paragraph{Paso 2 — Normalización por desviación estándar espacialmente promediada.} Para preservar los gradientes espaciales nativos de cada campo, la normalización utiliza un escalar global único por variable:

\begin{equation}
    Z_{V}(\tau, \mathbf{s}) = \frac{V'(\tau, \mathbf{s})}{\dfrac{1}{A}\displaystyle\int_{A} \sigma_{V}(\mathbf{s})\, d\mathbf{s}},
\end{equation}

donde $V'$ es el campo de anomalías, $\sigma_{V}(\mathbf{s})$ es la desviación estándar temporal en el punto $\mathbf{s}$ y $A$ es el área total del dominio. Este factor escalar preserva la estructura espacial de los gradientes mientras equilibra el peso de ambas variables en la matriz de covarianza.

\paragraph{Paso 3 — Concatenación espacial.} Los campos normalizados se concatenan para formar el vector de estado combinado:

\begin{equation}\label{eq:meof_vector}
    \mathbf{Y}(\tau) = \big[ Z_P(\tau),\; Z_W(\tau) \big].
\end{equation}

El análisis EOF aplicado a la matriz combinada $\mathbf{Y}(\tau)$ produce modos acoplados que explican la varianza conjunta del sistema. El MEOF$_1$ produce simultáneamente un patrón espacial para precipitación ($\text{EOF}_{1,P}$) y otro para viento ($\text{EOF}_{1,W}$), ambos extraídos del mismo modo dominante de covariabilidad conjunta.

\subsection{Construcción de los prototipos estructurales}\label{subsec:prototipos}
Para cumplir con el Objetivo Específico 4, se desarrolla un procedimiento de síntesis visual que integra en un único panel —para cada combinación región--fase— cuatro fuentes de información estadística derivadas de los análisis anteriores. Cada prototipo está compuesto por cuatro capas superpuestas en coordenadas lagrangianas $(\Delta x, \Delta y)$, dentro del dominio de $9.5^{\circ}$ de radio. La lógica del compositing centrado en el sistema, que hace posible esta superposición, fue presentada en la Sección~\ref{subsubsec:tb_compositing} de los Fundamentos.

Las dos primeras capas corresponden a las isolíneas de la función de densidad de probabilidad estimada (PDFe) de la ubicación de máximos de precipitación y viento. A partir de la superficie de densidad $\hat{f}(\mathbf{s})$ calculada según la Ecuación~\ref{eq:tb_kde_simple}, se extraen los contornos correspondientes a los percentiles 75 y 90 de la densidad acumulada (Ecuación~\ref{eq:nivel_contorno_kde}). Sobre cada superficie se señaliza mediante un marcador estelar la coordenada modal $(\Delta x_{\text{moda}}, \Delta y_{\text{moda}})$ definida por la Ecuación~\ref{eq:moda_kde}, indicando la posición donde estadísticamente es más frecuente encontrar el máximo de la variable.

Las capas tercera y cuarta muestran las isolíneas de los patrones espaciales del primer modo del análisis multivariado (MEOF$_1$), obtenidos según la Sección~\ref{subsec:eof_multivariado} para precipitación ($\text{EOF}_{1,P}$) y viento ($\text{EOF}_{1,W}$) respectivamente. Previo al trazado de contornos, se aplica un suavizado gaussiano ($\sigma = 2$) a cada campo. Los niveles de contorno se determinan mediante percentiles aplicados separadamente sobre los valores positivos (percentiles 33 y 66, representados con línea sólida) y negativos (percentiles 33 y 66 del valor absoluto, línea discontinua). Esta estrategia de umbralización relativa garantiza comparabilidad entre distintas combinaciones región--fase, independientemente de la magnitud absoluta del patrón espacial del MEOF. Las regiones de carga positiva indican sectores donde la variable tiende a presentar anomalías positivas respecto a la media cuando el modo dominante se activa; las de carga negativa, supresión relativa.

El resultado se organiza como una figura matricial de $3 \times 4$ paneles (3 regiones de ciclogénesis $\times$ 4 fases evolutivas), donde cada panel integra las cuatro capas descritas: PDFe de precipitación, PDFe de viento, $\text{EOF}_{1,P}$ de precipitación y $\text{EOF}_{1,W}$ de viento. La superposición permite identificar el grado de organización espacial entre ambos campos en el entorno del ciclón. Una alta coincidencia entre el núcleo del PDFe de precipitación y la región de carga positiva del EOF de viento indica una estructura frontal compacta; un desacoplamiento revela asimetrías en la distribución de los máximos, posiblemente asociadas a procesos locales de interacción océano-atmósfera o a la geometría del frente. Este análisis ofrece, además, un marco de referencia para detectar configuraciones que podrían favorecer eventos compuestos, aunque su objetivo principal es la caracterización morfológica de la fase.

La coordenada modal de cada variable, anotada en la leyenda interna de cada panel como $(\Delta x, \Delta y)$ en grados, cuantifica explícitamente el desplazamiento preferencial de los máximos respecto al centro del vórtice. Este procedimiento se ejecuta de forma independiente para la Población Global y el subconjunto p90, permitiendo la comparación directa de los patrones estructurales entre ambos estratos de intensidad.